\documentclass[10pt,a4paper]{amsart}
\usepackage[margin=19mm]{geometry}
\usepackage{amsmath,amssymb,mathtools}
\usepackage[scr=rsfs,cal=euler]{mathalfa}
\usepackage{xfrac}
\usepackage[unicode,hidelinks,bookmarks=false]{hyperref}
\newcommand{\ie}{i.e.\ }
\newcommand{\cf}{cf.\ }
\newcommand{\resp}{respectively\ }
\newcommand{\Subsection}{\subsection}
\providecommand{\nobreakdash}{\nobreak\hbox{-}}
\setlength{\emergencystretch}{2em}
\multlinegap0pt
\usepackage{float}
\newtheorem{proposition}{Proposition}
\title{Level-set physics-informed neural networks for domain inverse problems of gravimetry}
\author{Jingnan Yao, Wenbin Li, Jianliang Qian}
\date{Source: arXiv:2607.03772v1 (CC BY 4.0)}
\begin{document}
\maketitle

\noindent\emph{Source and licence.} Level-set physics-informed neural networks for domain inverse problems of gravimetry. Version arXiv:2607.03772v1 (CC BY 4.0), distributed by arXiv under the Creative Commons Attribution 4.0 International licence, http://creativecommons.org/licenses/by/4.0/, as recorded in the arXiv licence metadata for this exact version and shown on the abstract page https://arxiv.org/abs/2607.03772v1. Full licence notice: \texttt{licenses/arxiv-2607.03772-CC-BY-4.0.txt}.\par\smallskip

\noindent\textit{Editorial scope.} Counted unit: the article's proposition prop5 (source0.tex lines 392--394). The article's complete original proof of it is delivered with it (source0.tex lines 399--424, one \texttt{proof} environment of 26 lines, which the article places away from the statement). No mathematics is written, repaired or re-derived: the statement and the proof are the article's own lines, in the article's own notation, and no retained line is substituted or re-flowed.  Retained uncounted as the source-local context the proof needs: lines 83--85.  Nothing was omitted from any retained span, so the package carries no omission record.  No retained line contains a citation, so the wrapper carries no bibliography and no bibliography entry.  No retained line contains a figure environment, an image-inclusion command or drawing code, so no image is delivered and the package declares no asset dependency.  Editorial formatting only, in addition: an AMS \texttt{amsart} preamble that restates the article's own declarations and macros used by the retained lines, namely the environments \texttt{proposition} on separate counters, together with the standard packages those lines need.  The retained environments print in this document as Proposition 1 in order of appearance, whereas the article prints the same items with the numbering of its own full document.  The complete native source of the article as distributed in the arXiv e-print for this exact version is bundled unchanged as \texttt{sources/204417/source0.tex}.
\begin{equation}\label{eqn1}
U(\mathbf{y};\rho)=\gamma\int_\Omega K(\mathbf{y},\mathbf{x})\rho(\mathbf{x})\,\mathrm{d}\mathbf{x}\,.
\end{equation}

\begin{proposition} \label{prop5}
Let $\Omega_0$ be a bounded domain and $\Omega\subset\Omega_0$. Given that the density $\rho$ with $\operatorname{supp}\rho\subset\Omega$ satisfies $\rho\in L^\infty(\Omega)$, the gravity potential $U$ satisfies $U\in W^{1,\infty}(\Omega_0)$.
\end{proposition}

\begin{proof}
(Proposition \ref{prop5}) According to (\ref{eqn1}),
\[
|U(\mathbf{y})| \leq \gamma \|\rho\|_{L^\infty(\Omega)} \int_\Omega |K(\mathbf{y}, \mathbf{x})| \, \mathrm{d}\mathbf{x},\quad  |\nabla U(\mathbf{y})| \leq \gamma \|\rho\|_{L^\infty(\Omega)} \int_\Omega |\nabla_{\mathbf{y}}K(\mathbf{y}, \mathbf{x})| \, \mathrm{d}\mathbf{x}.
\]
Since $\Omega_0$ is bounded, there exists $r_0>\frac{1}{2}$ such that $\Omega\subset\Omega_0\subset B(\mathbf{0}, r_0)$ (a ball of radius $r_0$).

(1) In $\mathbf{R}^2$, $K(\mathbf{y},\mathbf{x})=-\frac{1}{2\pi}\mathrm{ln}|\mathbf{y}-\mathbf{x}|$ and $\nabla_{\mathbf{y}}K(\mathbf{y}, \mathbf{x})=-\frac{1}{2\pi}\frac{\mathbf{y}-\mathbf{x}}{|\mathbf{y}-\mathbf{x}|^2}$. It follows that $\forall\mathbf{y}\in\Omega_0$,
\[
|U(\mathbf{y})| \leq \gamma \|\rho\|_{L^\infty(\Omega)} \int_{|\mathbf{z}| \le 2 r_0} \frac{1}{2\pi}\big|\mathrm{ln}|\mathbf{z}|\big| \, \mathrm{d}\mathbf{z} = \gamma \|\rho\|_{L^\infty(\Omega)} \int_0^{2 r_0}|\mathrm{ln}r|r\,\mathrm{d}r\,.
\]
Since $2 r_0>1$, $\int_0^{2 r_0}|\mathrm{ln}r|r\,\mathrm{d}r=2r_0^2\mathrm{ln}(2r_0)-r_0^2+\frac{1}{2}$, which implies that
$
|U(\mathbf{y})| \leq \gamma \|\rho\|_{L^\infty(\Omega)}\left(2 r_0^2\mathrm{ln}(2 r_0)-r_0^2+\frac{1}{2}\right)\,.
$
Similarly, $\forall\mathbf{y}\in\Omega_0$,
\[
 |\nabla U(\mathbf{y})| \leq \gamma \|\rho\|_{L^\infty(\Omega)} \int_{|\mathbf{z}| \le 2 r_0}\frac{1}{2\pi |\mathbf{z}|} \, \mathrm{d}\mathbf{z}=\gamma \|\rho\|_{L^\infty(\Omega)} \int_0^{2 r_0}\mathrm{d}r=2r_0\gamma\|\rho\|_{L^\infty(\Omega)}\,.
\]

(2) In $\mathbf{R}^3$, $K(\mathbf{y},\mathbf{x})=\frac{1}{4\pi|\mathbf{y}-\mathbf{x}|}$ and $\nabla_{\mathbf{y}}K(\mathbf{y}, \mathbf{x})=-\frac{1}{4\pi}\frac{\mathbf{y}-\mathbf{x}}{|\mathbf{y}-\mathbf{x}|^3}$. It follows that $\forall\mathbf{y}\in\Omega_0$,
\begin{align*}
|U(\mathbf{y})| \leq \gamma \|\rho\|_{L^\infty(\Omega)} \int_{|\mathbf{z}| \le 2 r_0} \frac{1}{4\pi |\mathbf{z}|}\, \mathrm{d}\mathbf{z} = \gamma \|\rho\|_{L^\infty(\Omega)} \int_0^{2r_0}r\,\mathrm{d}r=2r_0^2 \gamma \|\rho\|_{L^\infty(\Omega)}\,, \\
|\nabla U(\mathbf{y})| \leq \gamma \|\rho\|_{L^\infty(\Omega)} \int_{|\mathbf{z}| \le 2r_0} \frac{1}{4\pi |\mathbf{z}|^2}\, \mathrm{d}\mathbf{z} = \gamma \|\rho\|_{L^\infty(\Omega)} \int_0^{2r_0} \mathrm{d}r=2r_0 \gamma \|\rho\|_{L^\infty(\Omega)}\,. 
\end{align*}
\end{proof}

\end{document}
